\documentclass[12pt]{article}
\usepackage[utf8]{inputenc}
\usepackage[T1]{fontenc}
\usepackage{hyperref}
\title{Mayawali Language Resources for K-12 Curriculum Development}
\author{Compiled by Claude Sonnet 5.5 (claude-sonnet-5-5) for Andr\'as Kornai}
\date{}
\begin{document}
\maketitle

\section{Introduction}
Mayawali (ISO 639-3: yxa; also Maiawali, Mayawarli, Majuli) is an extinct language of western Queensland, centred on the Diamantina Lakes area around Springvale and Cork Stations. According to the State Library of Queensland it is closely related to Pitta Pitta and Kunkalanya, and it has no known speakers and minimal documentation. The sources I found are few, and I found no teaching materials.

\section{Primary Educational Resources}
I found no primers, textbooks or school programmes for Mayawali.

\section{Linguistic and Reference Materials}
\begin{itemize}
\item State Library of Queensland, language week post on Maiawali, the main profile I found: \url{https://www.slq.qld.gov.au/blog/language-week-week-twenty-six-maiawali}
\item Wikipedia, Maiawali (people and language): \url{https://EN.wikipedia.com/wiki/Maiawali}
\item Breen, G. (1981), \emph{The Mayi Languages of the Queensland Gulf Country}, AIAS New Series 29, the standard comparative reference for the Mayi languages (library record): \url{https://catalog.hathitrust.org/Record/000106521/Cite}
\end{itemize}

\section{Religious and Community Resources}
I found no scripture or community language resources specific to Mayawali.

\section{Digital and Audio Resources}
I found no audio resources. The State Library of Queensland post above mentions earlier records: Breen recorded words from an informant at Boulia in 1969, and Blake recorded another speaker in 1972. Wordlists from pastoral properties also exist from the late 1800s.

\section{Curriculum Development Context}
Documentation is minimal, so any curriculum would be limited to a small word list and place names. It would need to be developed with descendant communities, and possibly with reference to the related Pitta Pitta language.
\end{document}